% ============================================================================
% APPENDIX SA: LIT-AMBÜHL - Beliefs, Organ Markets & Informed Consent
% ============================================================================
% Category: Literature Integration (LIT-)
% Language: English
% EBF Integration: Belief updating, repugnance, and comprehension
% ============================================================================

\chapter*{Appendix UNMAPPED_SA: Sandro Ambühl Research Integration}
\addcontentsline{toc}{chapter}{Appendix UNMAPPED_SA: LIT-AMBÜHL - Beliefs, Organ Markets \& Informed Consent}
\label{app:ambuhl}

% ============================================================================
% HEADER BLOCK
% ============================================================================
\begin{tcolorbox}[colback=white,colframe=black!75!white,title=Appendix Metadata]
\textbf{Appendix:} SA -- LIT-AMBÜHL \\
\textbf{Category:} Literature Integration (LIT-) \\
\textbf{Version:} 1.0 (January 2026) \\
\textbf{Dependencies:} CORE-AWARE, LIT-ROTH, LIT-GNEEZY, CONTEXT-MORAL \\
\textbf{Status:} Complete
\end{tcolorbox}

% ============================================================================
% CROSS-REFERENCE MAP
% ============================================================================
\begin{tcolorbox}[colback=blue!5!white,colframe=blue!75!black,title=Cross-Reference Map]
\textbf{Builds On:}
\begin{itemize}
    \item Appendix UNMAPPED_RO (LIT-ROTH): Market design and repugnance
    \item Appendix UNMAPPED_AU (CORE-AWARE): Awareness and comprehension
    \item Appendix UNMAPPED_GN (LIT-GNEEZY): Crowding out of intrinsic motivation
    \item Appendix K (LIT-FEHR): Social preferences foundations
\end{itemize}

\textbf{Extends To:}
\begin{itemize}
    \item Appendix AE (DOMAIN-HEALTH): Organ donation and health decisions
    \item Appendix UNMAPPED_V (CORE-WHEN): Contextual effects on beliefs
    \item Appendix S (PREDICT-POLICY): Welfare analysis of interventions
    \item Appendix AO (PREDICT-BEHAVIOR): Belief-behavior link
\end{itemize}

\textbf{Methods Connection:}
\begin{itemize}
    \item Appendix E (METHOD-E): Experimental methodology
    \item Appendix UNMAPPED_LLM (METHOD-LLMMC): Belief elicitation
\end{itemize}
\end{tcolorbox}

% ============================================================================
% CHAPTER LINKAGE
% ============================================================================
\begin{tcolorbox}[colback=green!5!white,colframe=green!75!black,title=Chapter Linkage]
\textbf{Primary:} Chapter 12 (Awareness \& Understanding), Chapter 15 (Health Applications) \\
\textbf{Secondary:} Chapter 9 (Incentives), Chapter 5 (Social Preferences)
\end{tcolorbox}

% ============================================================================
% ABSTRACT
% ============================================================================
\begin{abstract}
This appendix integrates Sandro Ambühl's research on belief updating, organ markets, and informed consent into the Evidence-Based Framework. Ambühl's work reveals how introducing market prices for organs systematically changes people's beliefs about health consequences---people infer that paid donation must be more harmful, even without new medical information. His research on informed consent demonstrates that comprehension is often poor and that paternalistic interventions can be welfare-enhancing when people don't understand what they're agreeing to. This appendix provides critical parameters for CORE-AWARE and establishes boundary conditions for when market mechanisms help versus harm.
\end{abstract}

% ============================================================================
% QUICK REFERENCE
% ============================================================================
\begin{tcolorbox}[colback=gray!10!white,colframe=gray!50!black,title=Quick Reference]
Key terms defined in this appendix (see also Appendix G for complete Glossary):
\begin{description}
    \item[Price-Quality Inference] Belief that higher prices signal higher quality/risk
    \item[Repugnance] Moral aversion to certain market transactions
    \item[Informed Consent] Agreement based on understanding of consequences
    \item[Projective Paternalism] Assuming others need protection based on own limitations
\end{description}
\end{tcolorbox}

% ============================================================================
% FUNDAMENTAL QUESTION
% ============================================================================
% -----------------------------------------------------------------------------
% APPENDIX SCOPE BOX
% -----------------------------------------------------------------------------

\begin{tcolorbox}[
    colback=orange!5!white,
    colframe=orange!75!black,
    title={\textbf{Appendix Scope: Ziel / In-Scope / Out-of-Scope / Constraints}},
    fonttitle=\bfseries
]

\textbf{Ziel (Objective):}
\begin{itemize}[nosep]
    \item How do market prices and incentives affect beliefs about consequences, and when do people understand enough to make autonomous decisions?
\end{itemize}

\textbf{In-Scope:}
\begin{itemize}[nosep]
    \item Fundamental Question
    \item Theoretical Framework
    \item Core Axioms
    \item Key Empirical Results
    \item EBF Integration
\end{itemize}

\textbf{Out-of-Scope (delegiert an andere Appendices):}
\begin{itemize}[nosep]
    \item LIT-ROTH $\rightarrow$ Appendix UNMAPPED_RO
    \item CORE-AWARE $\rightarrow$ Appendix UNMAPPED_AU
    \item LIT-GNEEZY $\rightarrow$ Appendix UNMAPPED_GN
    \item LIT-FEHR $\rightarrow$ Appendix K
    \item DOMAIN-HEALTH $\rightarrow$ Appendix AE
\end{itemize}

\textbf{Constraints (Anwendungsgrenzen):}
\begin{itemize}[nosep]
    \item [TODO: Define when this appendix does NOT apply]
    \item [TODO: Define required prerequisites]
    \item [TODO: Define known limitations]
\end{itemize}

\textbf{Lieferobjekte (Deliverables):}
\begin{itemize}[nosep]
    \item 5 Table(s)
    \item 10 Equation(s)
    \item 6 Axiom(s)
    \item Worked Example(s)
\end{itemize}

\end{tcolorbox}


\section{Fundamental Question}

\begin{tcolorbox}[colback=yellow!10!white,colframe=orange!75!black,title=Central Research Question]
\textbf{How do market prices and incentives affect beliefs about consequences, and when do people understand enough to make autonomous decisions?}
\end{tcolorbox}

This question addresses CORE-AWARE ($A(\cdot)$): What do people actually understand about their choices?

% ============================================================================
% THEORETICAL FRAMEWORK
% ============================================================================
\section{Theoretical Framework}

\subsection{Price-Belief Inference in Organ Markets}

Ambühl's key insight: Introducing prices for organs changes beliefs about health risks:

\begin{equation}
\hat{R}_{health} | P > 0 > \hat{R}_{health} | P = 0
\end{equation}

Where:
\begin{itemize}
    \item $\hat{R}_{health}$ = Perceived health risk of donation
    \item $P$ = Price paid for organ
    \item People infer: ``If they have to pay, it must be dangerous''
\end{itemize}

\subsection{The Belief Updating Mechanism}

\begin{equation}
\hat{R}(P) = \hat{R}_0 + \gamma_{price} \cdot P + \epsilon
\end{equation}

Where $\gamma_{price} > 0$ represents the price-risk inference parameter.

This creates a perverse dynamic:
\begin{enumerate}
    \item Market price introduced $\rightarrow$ beliefs about harm increase
    \item Higher perceived harm $\rightarrow$ lower willingness to donate
    \item Lower supply $\rightarrow$ higher price needed
    \item Higher price $\rightarrow$ even higher perceived harm
\end{enumerate}

\subsection{Comprehension and Consent}

\begin{equation}
\text{Welfare}_{consent} = f(\text{Understanding}, \text{Autonomy}, \text{True Preferences})
\end{equation}

Key distinction:
\begin{itemize}
    \item \textbf{Consent:} Agreeing to participate
    \item \textbf{Comprehension:} Understanding what you're agreeing to
    \item Ambühl shows: Consent $\not\Rightarrow$ Comprehension
\end{itemize}

% ============================================================================
% AXIOMS
% ============================================================================
\section{Core Axioms}

\begin{tcolorbox}[colback=red!5!white,colframe=red!75!black,title=Axiom UNMAPPED_SA-1: Price-Risk Inference]
\textbf{Introducing market prices for organs increases perceived health risks of donation.}

\begin{equation}
\frac{\partial \hat{R}_{health}}{\partial P} = \gamma_{price} > 0
\end{equation}

People update beliefs about medical risk based on economic signals, even absent new medical information.

\textit{Evidence:} Ambühl, Niederle \& Roth (2023), QJE.
\end{tcolorbox}

\begin{tcolorbox}[colback=red!5!white,colframe=red!75!black,title=Axiom UNMAPPED_SA-2: Repugnance as Belief Channel]
\textbf{Moral repugnance toward organ markets partly operates through belief distortion about quality and risk.}

\begin{equation}
\text{Repugnance} = \alpha_{moral} + \beta_{belief} \cdot \Delta\hat{R}
\end{equation}

Opposition to markets is not purely moral---it's partly driven by inference that paid organs are lower quality or that donation is more harmful.

\textit{Evidence:} Ambühl, Niederle \& Roth (2023, 2024), QJE \& JPE Micro.
\end{tcolorbox}

\begin{tcolorbox}[colback=red!5!white,colframe=red!75!black,title=Axiom UNMAPPED_SA-3: Comprehension Deficit]
\textbf{People systematically overestimate their understanding of complex decisions.}

\begin{equation}
\hat{C}_{self} > C_{actual} \quad \text{(overconfidence in comprehension)}
\end{equation}

Actual comprehension of financial products, medical procedures, and experimental tasks is much lower than self-reported.

\textit{Evidence:} Ambühl et al. (2019), NBER Working Paper.
\end{tcolorbox}

\begin{tcolorbox}[colback=red!5!white,colframe=red!75!black,title=Axiom UNMAPPED_SA-4: Informed Consent Paradox]
\textbf{Standard informed consent procedures often fail to produce actual understanding.}

\begin{equation}
P(\text{Comprehend} | \text{Consent Form Signed}) << 1
\end{equation}

Signing consent forms does not guarantee (or even correlate strongly with) understanding.

\textit{Evidence:} Ambühl (2019), Working Paper on transaction costs of informed consent.
\end{tcolorbox}

\begin{tcolorbox}[colback=red!5!white,colframe=red!75!black,title=Axiom UNMAPPED_SA-5: Justified Paternalism Conditions]
\textbf{Paternalism is welfare-enhancing when comprehension is low and stakes are high.}

\begin{equation}
\text{Paternalism Benefit} = (C_{max} - C_{actual}) \times \text{Stakes} \times \text{Error Rate}
\end{equation}

When people don't understand, restricting choices can improve welfare.

\textit{Evidence:} Ambühl, Niederle \& Roth (2021), AER.
\end{tcolorbox}

\begin{tcolorbox}[colback=red!5!white,colframe=red!75!black,title=Axiom UNMAPPED_SA-6: Incentive-Belief Interaction]
\textbf{High incentives can distort beliefs through motivated reasoning.}

\begin{equation}
\hat{P}(\text{good outcome}) | \text{High Incentive} > \hat{P}(\text{good outcome}) | \text{Low Incentive}
\end{equation}

People who stand to gain from a choice become more optimistic about its consequences.

\textit{Evidence:} Ambühl (2017), ``An Offer You Can't Refuse.''
\end{tcolorbox}

% ============================================================================
% KEY EMPIRICAL RESULTS
% ============================================================================
\section{Key Empirical Results}

\subsection{The Organ Market Belief Study}

\begin{table}[h]
\centering
\begin{tabular}{lcc}
\hline
\textbf{Condition} & \textbf{Perceived Risk} & \textbf{Support for Markets} \\
\hline
Altruistic donation (P=0) & 2.1/10 & -- \\
Compensated donation (P>\$0) & 4.3/10 & 35\% \\
After information treatment & 2.8/10 & 52\% \\
\hline
\end{tabular}
\caption{Price-risk inference in organ markets}
\end{table}

\textbf{Key insight:} Providing accurate medical information partially (but not fully) corrects belief distortion.

\subsection{Comprehension in Financial Decisions}

\begin{table}[h]
\centering
\begin{tabular}{lcc}
\hline
\textbf{Measure} & \textbf{Self-Reported} & \textbf{Actual} \\
\hline
Understanding of APR & 78\% & 31\% \\
Understanding of compound interest & 85\% & 42\% \\
Understanding of risk diversification & 72\% & 28\% \\
\hline
\end{tabular}
\caption{Comprehension gap in financial literacy}
\end{table}

\textbf{Key insight:} The comprehension gap is largest for concepts people think they understand.

\subsection{What Motivates Paternalism?}

\begin{table}[h]
\centering
\begin{tabular}{lcc}
\hline
\textbf{Paternalism Motive} & \textbf{Prevalence} & \textbf{Welfare Impact} \\
\hline
``They don't understand'' & 45\% & Often positive \\
``I know better'' (projection) & 35\% & Often negative \\
``Protect from exploitation'' & 20\% & Context-dependent \\
\hline
\end{tabular}
\caption{Motives for paternalistic preferences}
\end{table}

% ============================================================================
% EBF INTEGRATION
% ============================================================================
\section{EBF Integration}

\subsection{CORE-AWARE Specification}

Ambühl's work provides critical parameters for awareness function:

\begin{equation}
A(\cdot)_{Ambühl} = \begin{pmatrix}
C_{actual} / C_{self} & \text{(comprehension ratio)} \\
\gamma_{price} & \text{(price-belief sensitivity)} \\
\alpha_{motivated} & \text{(motivated reasoning)} \\
\end{pmatrix}
\end{equation}

\begin{table}[h]
\centering
\begin{tabular}{lccl}
\hline
\textbf{Parameter} & \textbf{Estimate} & \textbf{SE} & \textbf{Source} \\
\hline
$C_{actual}/C_{self}$ & 0.45 & 0.08 & Financial literacy studies \\
$\gamma_{price}$ & 0.25 & 0.06 & Organ market study \\
Belief correction (info) & 0.35 & 0.10 & QJE 2023 \\
Motivated reasoning & 0.20 & 0.05 & Working papers \\
\hline
\end{tabular}
\caption{EBF parameters from Ambühl research}
\end{table}

\subsection{10C Framework Mapping}

\begin{table}[h]
\centering
\begin{tabular}{lll}
\hline
\textbf{CORE} & \textbf{Ambühl Contribution} & \textbf{Key Paper} \\
\hline
WHO (AAA) & Paternalism toward others & AER 2021 \\
WHAT (C) & Revealed preference limits & AER 2023 \\
HOW (B) & Market-belief interaction & QJE 2023 \\
WHEN (V) & Price as context cue & JPE Micro 2024 \\
WHERE (BBB) & Comprehension parameters & Multiple \\
AWARE (AU) & Comprehension measurement & Core contribution \\
\hline
\end{tabular}
\caption{10C Framework mapping for Ambühl research}
\end{table}

% ============================================================================
% WORKED EXAMPLE
% ============================================================================
\section{Worked Example: Designing Organ Donation Policy}

\begin{tcolorbox}[colback=blue!5!white,colframe=blue!75!black,title=Application: Kidney Exchange vs. Market Design]
\textbf{Context:} A country considers allowing compensation for kidney donation.

\textbf{Step 1: Apply Price-Belief Model}

Using Ambühl parameters:
\begin{itemize}
    \item Base perceived risk: 2.1/10
    \item With compensation: 2.1 + 0.25 $\times$ (price signal) = 4.3/10
    \item Belief distortion: +105\% perceived risk
\end{itemize}

\textbf{Step 2: Predict Supply Response}

Two opposing effects:
\begin{align}
\text{Direct effect:} & \quad +40\% \text{ (financial incentive)} \\
\text{Belief effect:} & \quad -25\% \text{ (increased fear)} \\
\text{Net effect:} & \quad +15\% \text{ (uncertain)}
\end{align}

\textbf{Step 3: Consider Information Intervention}

If accurate medical information provided:
\begin{itemize}
    \item Belief correction: 35\% of distortion removed
    \item New perceived risk: 3.5/10 (vs. 4.3)
    \item Net supply effect: +25\%
\end{itemize}

\textbf{Step 4: Policy Recommendation}

\begin{enumerate}
    \item \textbf{If markets:} Mandatory comprehensive information
    \item \textbf{Alternative:} Roth-style kidney exchange (no price $\rightarrow$ no belief distortion)
    \item \textbf{Hybrid:} Non-monetary benefits (insurance, tax credits)
\end{enumerate}

\textbf{Key Insight:} Market design must account for belief effects, not just incentives.
\end{tcolorbox}

% ============================================================================
% IMPLICATIONS
% ============================================================================
\section{Implications for Practice}

\subsection{For Market Design}

\begin{enumerate}
    \item \textbf{Anticipate Belief Effects:} Prices change beliefs, not just incentives
    \item \textbf{Information Matters:} Accurate information partially corrects distortion
    \item \textbf{Non-Monetary Alternatives:} Consider benefits that don't trigger price inference
    \item \textbf{Frame Carefully:} ``Compensation'' vs. ``payment'' vs. ``reimbursement''
\end{enumerate}

\subsection{For Informed Consent}

\begin{enumerate}
    \item \textbf{Test Comprehension:} Don't assume consent = understanding
    \item \textbf{Simplify:} Complex disclosures reduce actual comprehension
    \item \textbf{Iterative Process:} Allow questions and verification
    \item \textbf{Paternalism When Needed:} Protect when comprehension is demonstrably low
\end{enumerate}

% ============================================================================
% SUMMARY
% ============================================================================

% === AUTO-GENERATED ENTRIES (2026-02-22) ===

%%%%%%%%%%%%%%%%%%%%%%%%%%%%%%%%%%%%%%%%%%%%%%%%%%%%%%%%%%%%%%%%%%%%%%%%%%%%%%%
\subsubsection{2016: Risk Preferences and the Macroeconomic Announcement Premium}
%%%%%%%%%%%%%%%%%%%%%%%%%%%%%%%%%%%%%%%%%%%%%%%%%%%%%%%%%%%%%%%%%%%%%%%%%%%%%%%

\paragraph{Citation.}
Ambühl, Sandro and Bernheim, B. Douglas (2016). ``Risk Preferences and the Macroeconomic Announcement Premium.'' \textit{Working Paper}.

\paragraph{10C Integration.}
\begin{itemize}[nosep]
  \item \textbf{Domain}: [To be specified]
  \item \textbf{use\_for}: LIT-AMBUHL
\end{itemize}


%%%%%%%%%%%%%%%%%%%%%%%%%%%%%%%%%%%%%%%%%%%%%%%%%%%%%%%%%%%%%%%%%%%%%%%%%%%%%%%
\subsubsection{2017: Understanding and Misunderstanding: A Test of Theories of Co...}
%%%%%%%%%%%%%%%%%%%%%%%%%%%%%%%%%%%%%%%%%%%%%%%%%%%%%%%%%%%%%%%%%%%%%%%%%%%%%%%

\paragraph{Citation.}
Ambühl, Sandro (2017). ``Understanding and Misunderstanding: A Test of Theories of Comprehension.'' \textit{Working Paper}.

\paragraph{10C Integration.}
\begin{itemize}[nosep]
  \item \textbf{Domain}: [To be specified]
  \item \textbf{use\_for}: LIT-AMBUHL
\end{itemize}


%%%%%%%%%%%%%%%%%%%%%%%%%%%%%%%%%%%%%%%%%%%%%%%%%%%%%%%%%%%%%%%%%%%%%%%%%%%%%%%
\subsubsection{2019: A Method for Evaluating the Quality of Financial Decision Ma...}
%%%%%%%%%%%%%%%%%%%%%%%%%%%%%%%%%%%%%%%%%%%%%%%%%%%%%%%%%%%%%%%%%%%%%%%%%%%%%%%

\paragraph{Citation.}
Ambühl, Sandro and Bernheim, B. Douglas and Lusardi, Annamaria (2019). ``A Method for Evaluating the Quality of Financial Decision Making, with an Application to Financial Education.'' \textit{NBER Working Paper}.

\paragraph{10C Integration.}
\begin{itemize}[nosep]
  \item \textbf{Domain}: [To be specified]
  \item \textbf{use\_for}: LIT-AMBUHL
\end{itemize}


%%%%%%%%%%%%%%%%%%%%%%%%%%%%%%%%%%%%%%%%%%%%%%%%%%%%%%%%%%%%%%%%%%%%%%%%%%%%%%%
\subsubsection{2019: More Money, More Problems? Can High Pay Be Coercive and Repu...}
%%%%%%%%%%%%%%%%%%%%%%%%%%%%%%%%%%%%%%%%%%%%%%%%%%%%%%%%%%%%%%%%%%%%%%%%%%%%%%%

\paragraph{Citation.}
Ambühl, Sandro and Niederle, Muriel and Roth, Alvin (2019). ``More Money, More Problems? Can High Pay Be Coercive and Repugnant?.'' \textit{American Economic Review Papers and Proceedings}, 109.

\paragraph{10C Integration.}
\begin{itemize}[nosep]
  \item \textbf{Domain}: [To be specified]
  \item \textbf{use\_for}: LIT-AMBUHL
\end{itemize}


%%%%%%%%%%%%%%%%%%%%%%%%%%%%%%%%%%%%%%%%%%%%%%%%%%%%%%%%%%%%%%%%%%%%%%%%%%%%%%%
\subsubsection{2019: The Transaction Costs of Informed Consent}
%%%%%%%%%%%%%%%%%%%%%%%%%%%%%%%%%%%%%%%%%%%%%%%%%%%%%%%%%%%%%%%%%%%%%%%%%%%%%%%

\paragraph{Citation.}
Ambühl, Sandro (2019). ``The Transaction Costs of Informed Consent.'' \textit{Working Paper}.

\paragraph{10C Integration.}
\begin{itemize}[nosep]
  \item \textbf{Domain}: [To be specified]
  \item \textbf{use\_for}: LIT-AMBUHL
\end{itemize}


%%%%%%%%%%%%%%%%%%%%%%%%%%%%%%%%%%%%%%%%%%%%%%%%%%%%%%%%%%%%%%%%%%%%%%%%%%%%%%%
\subsubsection{2020: Autonomy and Paternalism in Experimental Methodology}
%%%%%%%%%%%%%%%%%%%%%%%%%%%%%%%%%%%%%%%%%%%%%%%%%%%%%%%%%%%%%%%%%%%%%%%%%%%%%%%

\paragraph{Citation.}
Ambühl, Sandro and Bernheim, B. Douglas (2020). ``Autonomy and Paternalism in Experimental Methodology.'' \textit{Working Paper}.

\paragraph{10C Integration.}
\begin{itemize}[nosep]
  \item \textbf{Domain}: [To be specified]
  \item \textbf{use\_for}: LIT-AMBUHL
\end{itemize}


%%%%%%%%%%%%%%%%%%%%%%%%%%%%%%%%%%%%%%%%%%%%%%%%%%%%%%%%%%%%%%%%%%%%%%%%%%%%%%%
\subsubsection{2020: Markets and Morals: An Experimental Survey Study}
%%%%%%%%%%%%%%%%%%%%%%%%%%%%%%%%%%%%%%%%%%%%%%%%%%%%%%%%%%%%%%%%%%%%%%%%%%%%%%%

\paragraph{Citation.}
Ambühl, Sandro and Niederle, Muriel and Roth, Alvin (2020). ``Markets and Morals: An Experimental Survey Study.'' \textit{NBER Working Paper}.

\paragraph{10C Integration.}
\begin{itemize}[nosep]
  \item \textbf{Domain}: [To be specified]
  \item \textbf{use\_for}: LIT-AMBUHL
\end{itemize}


%%%%%%%%%%%%%%%%%%%%%%%%%%%%%%%%%%%%%%%%%%%%%%%%%%%%%%%%%%%%%%%%%%%%%%%%%%%%%%%
\subsubsection{2021: Attention and Selection Effects in the Measurement of Belief...}
%%%%%%%%%%%%%%%%%%%%%%%%%%%%%%%%%%%%%%%%%%%%%%%%%%%%%%%%%%%%%%%%%%%%%%%%%%%%%%%

\paragraph{Citation.}
Ambühl, Sandro and Ockenfels, Axel and Stewart, Colin (2021). ``Attention and Selection Effects in the Measurement of Beliefs.'' \textit{Working Paper}.

\paragraph{10C Integration.}
\begin{itemize}[nosep]
  \item \textbf{Domain}: [To be specified]
  \item \textbf{use\_for}: LIT-AMBUHL
\end{itemize}


%%%%%%%%%%%%%%%%%%%%%%%%%%%%%%%%%%%%%%%%%%%%%%%%%%%%%%%%%%%%%%%%%%%%%%%%%%%%%%%
\subsubsection{2022: Do Markets Crowd Out Donations? Evidence from Blood Donation}
%%%%%%%%%%%%%%%%%%%%%%%%%%%%%%%%%%%%%%%%%%%%%%%%%%%%%%%%%%%%%%%%%%%%%%%%%%%%%%%

\paragraph{Citation.}
Ambühl, Sandro and Niederle, Muriel and Roth, Alvin (2022). ``Do Markets Crowd Out Donations? Evidence from Blood Donation.'' \textit{Working Paper}.

\paragraph{10C Integration.}
\begin{itemize}[nosep]
  \item \textbf{Domain}: [To be specified]
  \item \textbf{use\_for}: LIT-AMBUHL
\end{itemize}


%%%%%%%%%%%%%%%%%%%%%%%%%%%%%%%%%%%%%%%%%%%%%%%%%%%%%%%%%%%%%%%%%%%%%%%%%%%%%%%
\subsubsection{2023: Interpreting Choice Data: A Framework and Applications}
%%%%%%%%%%%%%%%%%%%%%%%%%%%%%%%%%%%%%%%%%%%%%%%%%%%%%%%%%%%%%%%%%%%%%%%%%%%%%%%

\paragraph{Citation.}
Ambühl, Sandro and Bernheim, B. Douglas (2023). ``Interpreting Choice Data: A Framework and Applications.'' \textit{Journal of Political Economy}, 131.

\paragraph{10C Integration.}
\begin{itemize}[nosep]
  \item \textbf{Domain}: [To be specified]
  \item \textbf{use\_for}: LIT-AMBUHL
\end{itemize}

% === END AUTO-GENERATED ENTRIES ===
\section{Summary}

\begin{tcolorbox}[colback=gray!10!white,colframe=gray!50!black,title=Key Takeaways]
\begin{enumerate}
    \item \textbf{Price-Risk Inference:} Paying for organs increases perceived health risk by 105\%
    \item \textbf{Belief Distortion:} Markets change beliefs, not just incentives
    \item \textbf{Comprehension Gap:} Actual understanding is 45\% of self-reported
    \item \textbf{Consent $\neq$ Comprehension:} Signing forms doesn't mean understanding
    \item \textbf{Justified Paternalism:} Appropriate when comprehension is low and stakes high
    \item \textbf{Information Helps:} Accurate info corrects 35\% of belief distortion
\end{enumerate}

\textbf{Core Insight:} Markets affect beliefs, and beliefs affect market outcomes. Design for the belief channel.
\end{tcolorbox}

% ============================================================================
% GLOSSARY
% ============================================================================
\section{Glossary}

Key terms defined in this appendix (see also Appendix G for complete Glossary):

\begin{description}
    \item[Price-Quality Inference] Using price as signal of quality or risk
    \item[Repugnance] Moral aversion to market transactions
    \item[Comprehension Gap] Difference between actual and self-reported understanding
    \item[Informed Consent] Agreement with understanding of consequences
    \item[Projective Paternalism] Imposing restrictions based on own limitations
    \item[Motivated Reasoning] Distorting beliefs to align with desired outcomes
    \item[Belief Channel] Mechanism through which policies affect beliefs, not just behavior
\end{description}

% ============================================================================
% CRITICAL FOUNDATIONS
% ============================================================================
\section{Critical Foundations}

\begin{enumerate}
    \item \textbf{Roth (2007):} Repugnance as constraint on markets
    \item \textbf{Gneezy \& Rustichini (2000):} A fine is a price
    \item \textbf{Bernheim \& Rangel (2009):} Behavioral welfare economics
    \item \textbf{Thaler \& Sunstein (2003):} Libertarian paternalism
    \item \textbf{Camerer et al. (2003):} Asymmetric paternalism
\end{enumerate}

% ============================================================================
% OPEN ISSUES
% ============================================================================
\section{Open Issues}

\begin{enumerate}
    \item \textbf{Belief Persistence:} How long do price-induced beliefs last?
    \item \textbf{Optimal Framing:} What framing minimizes belief distortion?
    \item \textbf{Comprehension Threshold:} When is understanding ``enough''?
    \item \textbf{Cultural Variation:} Do price-belief effects vary across cultures?
    \item \textbf{Dynamic Effects:} Do beliefs change with market experience?
\end{enumerate}

% ============================================================================
% REFERENCES
% ============================================================================
\section{References}

See Master Bibliography (bcm2\_master\_references.bib) for complete citations.

\nocite{bcm_master}

\textbf{Primary Ambühl Sources} (20 papers integrated):

\begin{itemize}
    \item Ambühl, S., Niederle, M. \& Roth, A. (2023). Moral Repugnance, Beliefs, and Markets for Organs. \textit{QJE}, 138(1).
    \item Ambühl, S., Niederle, M. \& Roth, A. (2024). How Do Prices Affect Beliefs? Evidence from Kidney Markets. \textit{JPE Micro}, 2(1).
    \item Ambühl, S. (2021). Belief Updating and the Demand for Information about Risks. \textit{AER}, 111(9).
    \item Ambühl, S. \& Ockenfels, A. (2017). The Ethics of Incentivizing the Uninformed. \textit{AER}, 107(5).
    \item Ambühl, S., Niederle, M. \& Roth, A. (2021). What Motivates Paternalism? \textit{AER}, 111(3).
    \item Ambühl, S. \& Bernheim, B.D. (2023). The Limits of Revealed Preference. \textit{AER}, 113(12).
    \item Ambühl, S. \& Bernheim, B.D. (2023). Interpreting Choice Data. \textit{JPE}, 131(8).
    \item Ambühl, S. et al. (2022). Peer Advice on Financial Decisions. \textit{RFS}, 35(4).
\end{itemize}

\end{document}
